%% toc: P3 40 40.8
\chapter{Nonconvergence and the Limits of Parity Definability}
\label{sec:p03-w03-p03-071-nonconvergence}

This chapter treats two negative results about first-order sentences on the
random ordered graph \(G_<(n,p)\), in which the vertices carry a built-in
total order.  The first is that the zero--one law and even mere convergence
fail: there is a first-order sentence \(A\) for which
\(\lim_{n\to\infty}\Pr[G_<(n,\frac12)\entails A]\) does not exist.  The proof
encodes arithmetic on an initial segment \(\{1,\ldots,u\}\) of the order
through the neighbourhoods of further vertices, measures how long an initial
segment can be encoded, and lets the sentence ask a question about the
logarithm of that length.  The second result runs the other way: for every
first-order \(A\) the difference
\(\Pr[G_<(n+1,\frac12)\entails A]-\Pr[G_<(n,\frac12)\entails A]\) tends to
\(0\).  Probabilities may fail to converge, but they cannot alternate from
one value of \(n\) to the next, and so no first-order sentence almost surely
expresses that the number of vertices is even.

The inputs are a first-order sentence in the language with the order symbol
\(<\), the adjacency symbol \(\sim\) and equality, and the edge probability
\(p=\frac12\).  The conclusions are the existence of a sentence whose
probabilities have \(\liminf 0\) and \(\limsup 1\), an explicit such
sentence, the almost sure representation of every \(k\)-ary relation on a
short initial segment of the order, first-order definitions of primality
and of Fermat primality for the length of an encoded segment, a
total-influence estimate for bounded-depth circuits with a complete proof
(including the switching lemma it rests on), the slow-oscillation bound for
every sentence with a rate, and the impossibility of an almost sure
first-order representation of evenness.

The presentation follows Spencer's author version of ``Logic and Random
Structures'' \cite{Spencer2007LogicRandomStructuresAuthor}, sections 12 and
13 (pp.~22--26).  That text is formatted separately from the chapter printed
in the 2007 Springer volume \emph{Finite Model Theory and Its Applications};
its sections are numbered without a chapter prefix, and all locators below
are to it.  Spencer gives both results in outline.  Three points of the
outline do not hold as written, and the chapter corrects them explicitly:
the literal neighbourhood encoding cannot represent addition
(\cref{thm:p03-w03-p03-071-literal-fails}); the upper limit on the encodable
length needs a restriction on the witness interval
(\cref{def:p03-w03-p03-071-clean}); and the circuit lower bound the source
names, that bounded-depth circuits cannot compute majority, does not by
itself give the needed estimate, so the proof uses a total-influence bound
instead, proved in full here (\cref{thm:p03-w03-p03-071-ac0-influence}).

\section{The random ordered graph}
\label{sec:p03-w03-p03-071-ordered}

\begin{hbconvention}{def:p03-w03-p03-071-assumptions}{Standing assumptions}
The language has the binary symbols \(<\) and \(\sim\) and equality.
Formulas are built from atomic formulas with \(\neg,\wedge,\vee,\exists,
\forall\); \(\to\) and \(\leftrightarrow\) abbreviate their usual
definitions.  For a first-order sentence \(A\) and a random structure on
\(n\) points, \(A\) holds \emph{almost surely} if \(\Pr[A]\to1\) and
\emph{almost never} if \(\Pr[A]\to0\) as \(n\to\infty\)
\cite[p.~2]{Spencer2007LogicRandomStructuresAuthor}.  The logarithm
\(\log\) is to base \(2\), as in the source
\cite[p.~23]{Spencer2007LogicRandomStructuresAuthor}, and
\(L=L(n)=(\log n)^{1/3}\).  For a real \(t\), \(\lfloor t\rfloor\) is the
largest integer not above \(t\).
\end{hbconvention}

\begin{hbdefinition}{def:p03-w03-p03-071-gord}{The random ordered graph \(G_<(n,p)\) and \(f_A(n)\)}
Let \(n\geq1\) and \(0\leq p\leq1\).  The random ordered graph \(G_<(n,p)\)
has vertex set \(\Omega=\{1,\ldots,n\}\), totally ordered by the built-in
relation \(<\), which is the usual order of the integers and belongs to the
language.  Adjacency \(\sim\) is irreflexive and symmetric, and each of the
\(\binom n2\) pairs is adjacent with probability \(p\), independently of the
others \cite[p.~22]{Spencer2007LogicRandomStructuresAuthor}.  For a sentence
\(A\),
\[
  f_A(n)=\Pr[G_<(n,p)\entails A]
\]
\cite[p.~25]{Spencer2007LogicRandomStructuresAuthor}.  The element \(1\) is
first-order definable as the element with nothing below it, and
\(y=x+1\) abbreviates \(x<y\wedge\neg\exists z\,(x<z\wedge z<y)\); when
\(y\neq1\), \(x=y-1\) means \(y=x+1\).  Addition and subtraction are not in
the language \cite[p.~22]{Spencer2007LogicRandomStructuresAuthor}.  For
\(u\in\Omega\) and a vertex \(e\), the \emph{trace} of \(e\) on
\(\{1,\ldots,u\}\) is \(N_u(e)=\{w\leq u: w\sim e\}\).
\end{hbdefinition}

\begin{hbtheorem}{thm:p03-w03-p03-071-no-zero-one}{proved=annex:p03-w03-p03-071-no-zero-one}{No zero--one law with a built-in order}
Let \(B\) be the sentence
\[
  \exists_x\exists_y\,(x\neq y)\wedge(x<y)\wedge[\forall_z\,z<y\to z=x]
  \wedge x\sim y
\]
\cite[p.~22]{Spencer2007LogicRandomStructuresAuthor}, which says that the
first two elements are adjacent.  For every \(n\geq2\),
\(\Pr[G_<(n,p)\entails B]=p\).  In particular for \(p=\frac12\) the
sentence \(B\) holds neither almost surely nor almost never.
\end{hbtheorem}

\section{Arithmetization and its almost sure encoding}
\label{sec:p03-w03-p03-071-arithmetization}

\begin{hbdefinition}{def:p03-w03-p03-071-arithmetizable}{Second-order arithmetization of \(\{1,\ldots,u\}\)}
Let \(S=\{1,\ldots,u\}\) with the induced order.  Ternary relations
\(+,\ast\subseteq S^3\), read as \(x+y=z\) and \(x\cdot y=z\), \emph{arithmetize}
\(S\) if for all \(x,y,z\in S\):
\begin{enumerate}
\item \(+(x,1,z)\) if and only if \(z=x+1\);
\item when \(y\neq1\), \(+(x,y,z)\) if and only if \(z\neq1\) and
  \(+(x,y-1,z-1)\);
\item \(\ast(x,1,z)\) if and only if \(z=x\);
\item when \(y\neq1\), \(\ast(x,y,z)\) if and only if there exists \(v\in S\)
  with \(\ast(x,y-1,v)\) and \(+(x,v,z)\).
\end{enumerate}
\(S\) is \emph{arithmetizable} if such relations exist
\cite[p.~22]{Spencer2007LogicRandomStructuresAuthor}.  The clause
``\(z\neq1\)'' in item~2 makes explicit that \(z-1\) must exist; the source
leaves it implicit.  Write \(\mathrm{Add}_u=\{(x,y,z)\in S^3:x+y=z\}\) and
\(\mathrm{Mul}_u=\{(x,y,z)\in S^3:xy=z\}\), with the arithmetic of the
integers.
\end{hbdefinition}

\begin{hblemma}{thm:p03-w03-p03-071-unique}{proved=annex:p03-w03-p03-071-unique}{The arithmetization is unique}
Relations \(+,\ast\subseteq S^3\) arithmetize \(S=\{1,\ldots,u\}\) if and only
if \(+=\mathrm{Add}_u\) and \(\ast=\mathrm{Mul}_u\).  In particular every
\(\{1,\ldots,u\}\) is arithmetizable, by exactly one pair of relations.
\end{hblemma}

\begin{hbdefinition}{def:p03-w03-p03-071-literal}{Spencer's encoding \(R_{c,d}\)}
For \(c\leq d\) in \(\Omega\) and \(S=\{1,\ldots,u\}\), Spencer writes
\(R_{c,d}(x,y,z)\) if \(x,y\leq z\) and there is \(e\) with \(c\leq e<d\)
adjacent to \(x\), \(y\) and \(z\) and to no other element of \(S\); and he
calls \(S\) \emph{first order arithmetizable} if there are \(c,d\) and
\(c',d'\) such that \(R_{c,d}\) and \(R_{c',d'}\) have the four properties
of \(+\) and \(\ast\) in \cref{def:p03-w03-p03-071-arithmetizable}
\cite[pp.~22--23]{Spencer2007LogicRandomStructuresAuthor}.  The relation
\(R_{c,d}\) is first-order definable from the parameters \(c,d,u\), because
``\(w\in S\)'' is \(w\leq u\).
\end{hbdefinition}

\begin{hbtheorem}{thm:p03-w03-p03-071-literal-fails}{proved=annex:p03-w03-p03-071-literal-fails}{The literal encoding never represents addition}
For every ordered graph on \(\Omega\), every \(u\geq2\) in \(\Omega\) and
all \(c\leq d\) in \(\Omega\), \(R_{c,d}\neq\mathrm{Add}_u\) and
\(R_{c,d}\neq\mathrm{Mul}_u\).  Consequently no \(\{1,\ldots,u\}\) with
\(u\geq2\) is first order arithmetizable in the sense of
\cref{def:p03-w03-p03-071-literal}, and Spencer's representation lemma
\cite[p.~23]{Spencer2007LogicRandomStructuresAuthor} fails as literally
stated for every \(u\geq2\).
\end{hbtheorem}

The defect is that a vertex whose trace is a two-element set
\(\{a,b\}\) with \(a<b\) witnesses three triples, \((a,a,b)\), \((a,b,b)\) and
\((b,a,b)\), and no arithmetic relation contains all three.  The repair below
keeps Spencer's idea, a vertex per triple whose trace on \(S\) is exactly
the triple, and changes three details: the inequalities become strict, the
unit cases of multiplication are written out, and every vertex of a witness
interval must have a trace of size \(2\) or \(3\).  The last clause is what
makes the upper bound of \cref{thm:p03-w03-p03-071-upper} provable.

\begin{hbdefinition}{def:p03-w03-p03-071-clean}{The corrected encoding and the formula \(\operatorname{Ar}(u)\)}
All quantifiers over elements of \(S=\{1,\ldots,u\}\) below are relativized
to \(w\leq u\).  For parameters \(u,c,d\) put
\begin{align*}
  \operatorname{Tr}_u(e;x,y,z)&:\quad x<z\wedge y<z\wedge e\sim x\wedge e\sim y
     \wedge e\sim z\\
  &\qquad\wedge\forall_{w\leq u}\,(e\sim w\to w=x\vee w=y\vee w=z),\\
  R^<_{c,d}(x,y,z)&:\quad \exists_e\,(c\leq e<d\wedge
     \operatorname{Tr}_u(e;x,y,z)),\\
  \operatorname{Clean}_u(c,d)&:\quad c\leq d\wedge u<c\\
  &\qquad\wedge\forall_e\,\bigl(c\leq e<d\to\exists_{x,y,z\leq u}
     \operatorname{Tr}_u(e;x,y,z)\bigr).
\end{align*}
Thus \(\operatorname{Tr}_u(e;x,y,z)\) says \(N_u(e)=\{x,y,z\}\) with
\(x,y<z\), and \(\operatorname{Clean}_u(c,d)\) says the interval
\([c,d)=\{e:c\leq e<d\}\) lies above \(S\) and each of its vertices has a
trace on \(S\) of exactly two or three elements.  With
\begin{align*}
  P(x,y,z)&:\quad R^<_{c,d}(x,y,z),\\
  T(x,y,z)&:\quad (y=1\wedge z=x)\vee(x=1\wedge z=y)\vee R^<_{c',d'}(x,y,z),
\end{align*}
let \(\operatorname{Good}(u,c,d,c',d')\) be the conjunction of
\(\operatorname{Clean}_u(c,d)\), \(\operatorname{Clean}_u(c',d')\) and the
four properties of \cref{def:p03-w03-p03-071-arithmetizable} for \(P\) in
place of \(+\) and \(T\) in place of \(\ast\).  Finally
\[
  \operatorname{Ar}(u):\quad\exists_c\exists_d\exists_{c'}\exists_{d'}\,
  \operatorname{Good}(u,c,d,c',d').
\]
The witness families are
\(\mathcal F_+(u)=\bigl\{\{a,b,a+b\}:1\leq a\leq b,\ a+b\leq u\bigr\}\) and
\(\mathcal F_\times(u)=\bigl\{\{a,b,ab\}:2\leq a\leq b,\ ab\leq u\bigr\}\), of
sizes \(m_+(u)\) and \(m_\times(u)\).
\end{hbdefinition}

\begin{hblemma}{thm:p03-w03-p03-071-witness}{proved=annex:p03-w03-p03-071-witness}{What a witness interval is}
Fix an ordered graph on \(\Omega\) and \(u\in\Omega\).
\begin{enumerate}
\item If \(u<c\), \(d=c+m_+(u)\leq n+1\) and the traces
  \(N_u(c),\ldots,N_u(d-1)\) are the sets of \(\mathcal F_+(u)\), each once,
  then \(\operatorname{Clean}_u(c,d)\) holds and \(R^<_{c,d}=\mathrm{Add}_u\)
  on \(S\).  If the same holds with \(\mathcal F_\times(u)\) and
  \(d'=c'+m_\times(u)\), then \(\operatorname{Clean}_u(c',d')\) holds and
  \(T=\mathrm{Mul}_u\).
\item \(\operatorname{Good}(u,c,d,c',d')\) holds if and only if
  \(\operatorname{Clean}_u(c,d)\), \(\operatorname{Clean}_u(c',d')\),
  \(P=\mathrm{Add}_u\) and \(T=\mathrm{Mul}_u\).
\item If \(\operatorname{Clean}_u(c,d)\) and \(P=\mathrm{Add}_u\), then each
  set of \(\mathcal F_+(u)\) is the trace of a vertex of \([c,d)\), so
  \(d-c\geq m_+(u)\), and the vertices \(c,\ldots,c+m_+(u)-1\) lie in
  \((u,n]\) and have traces of size \(2\) or \(3\).
\end{enumerate}
\end{hblemma}

\begin{hbtheorem}{thm:p03-w03-p03-071-representation}{proved=annex:p03-w03-p03-071-representation}{Representation lemma}
In \(G_<(n,\frac12)\), with probability tending to \(1\),
\(\operatorname{Ar}(u)\) holds for every \(u\) with \(1\leq u\leq0.9L\).
\end{hbtheorem}

This is Spencer's representation lemma with his bound
\(u\leq0.9\log^{1/3}n\) \cite[p.~23]{Spencer2007LogicRandomStructuresAuthor},
for the corrected encoding and uniformly in \(u\).  The proof is his: a block
of \(s\) consecutive vertices above \(S\) carries the required pattern with
probability \(2^{-us}\), and disjoint blocks are independent.

Spencer adds that with further technical work, perhaps with a different
bound on \(u\), every ternary and even every \(k\)-ary relation could almost
surely be represented \cite[p.~23]{Spencer2007LogicRandomStructuresAuthor}.
He specifies neither the encoding nor the bound.  The encoding
\(R^<_{c,d}\) of \cref{def:p03-w03-p03-071-clean} cannot do this: a trace is
a set, so \(R^<_{c,d}(x,y,z)\) holds exactly when \(R^<_{c,d}(y,x,z)\) does,
and only relations with \(x,y<z\) occur.  The block encoding below reads a tuple from \(k+1\)
consecutive vertices, uses the first as a marker, and so represents
arbitrary relations; the price is a smaller range of \(u\).

\begin{hbdefinition}{def:p03-w03-p03-071-block-encoding}{Block encoding of a \(k\)-ary relation}
Let \(k\geq1\).  For parameters \(u,c,d\) and variables
\(x_1,\ldots,x_k\), let \(\operatorname{Rep}^{(k)}_{u,c,d}(x_1,\ldots,x_k)\)
abbreviate the formula
\begin{align*}
  &\exists_{e_0}\cdots\exists_{e_k}\Bigl[c\leq e_0\wedge e_k<d\\
  &\qquad\wedge\bigwedge_{m<k}e_{m+1}=e_m+1
  \wedge\forall_{w\leq u}\,\neg(e_0\sim w)\\
  &\qquad\wedge\bigwedge_{m=1}^{k}\bigl(e_m\sim x_m\wedge
  \forall_{w\leq u}(e_m\sim w\to w=x_m)\bigr)\Bigr].
\end{align*}
On \(S=\{1,\ldots,u\}\) this formula defines the set of tuples
\((x_1,\ldots,x_k)\in S^k\) for which some \(k+1\) consecutive vertices
\(e_0<\cdots<e_k\) of the interval \([c,d)\) have traces
\(N_u(e_0)=\varnothing\) and \(N_u(e_m)=\{x_m\}\) for \(1\leq m\leq k\).  A
relation \(R\subseteq S^k\) is \emph{represented by} \((c,d)\) if
\(u<c\leq d\) and this set equals \(R\).  Put
\[
  U_k(n)=\Bigl\lfloor\Bigl(\frac{\log n}{2(k+1)}\Bigr)^{1/(k+1)}\Bigr\rfloor .
\]
\end{hbdefinition}

\begin{hbtheorem}{thm:p03-w03-p03-071-kary}{proved=annex:p03-w03-p03-071-kary}{Every \(k\)-ary relation on a short initial segment is represented}
Fix \(k\geq1\).  In \(G_<(n,\frac12)\), with probability tending to \(1\),
for every \(u\) with \(1\leq u\leq U_k(n)\) and every relation
\(R\subseteq\{1,\ldots,u\}^k\) there are \(c,d\in\Omega\) with \(u<c\leq d\)
such that \(R\) is represented by \((c,d)\).
\end{hbtheorem}

On the event of \cref{thm:p03-w03-p03-071-kary}, as \((c,d)\) ranges over
the pairs with \(u<c\leq d\), the formula
\(\operatorname{Rep}^{(k)}_{u,c,d}\) ranges over exactly the \(k\)-ary
relations on \(\{1,\ldots,u\}\): every pair defines some relation, and every
relation is defined by some pair.  A first-order quantifier over \(c,d\) thus
acts, almost surely, as a second-order quantifier over all \(k\)-ary
relations on \(\{1,\ldots,u\}\) for \(u\leq U_k(n)\).  This is the exact form
in which Spencer's remark is established here.  For \(k=3\) the range is
\(u\leq((\log n)/8)^{1/4}\), of smaller order than the \(0.9(\log n)^{1/3}\) of
\cref{thm:p03-w03-p03-071-representation}, which concerns only addition and
multiplication.  Whether every ternary relation is almost surely
representable, by this or another first-order encoding, for all \(u\) up to
a constant multiple of \((\log n)^{1/3}\) is not established in this chapter.

\begin{hbtheorem}{thm:p03-w03-p03-071-upper}{proved=annex:p03-w03-p03-071-upper}{Upper limit of the encoding}
In \(G_<(n,\frac12)\), for every \(u\geq30\),
\(\Pr[\operatorname{Ar}(u)]\leq n\,2^{-u^3/10}\).  Consequently, with
probability tending to \(1\), \(\operatorname{Ar}(u)\) fails for every
\(u\in\Omega\) with \(u\geq3L\).
\end{hbtheorem}

Spencer states that the representation fails almost surely above
\(C\log^{1/3}n\) for a computable constant \(C\), takes \(C=900\), and
excludes the argument \cite[p.~23]{Spencer2007LogicRandomStructuresAuthor}.
The bound proved here has \(C=3\).  Its proof needs the cleanliness clause:
without it, a witness interval may be arbitrarily long, and the count of
\cref{thm:p03-w03-p03-071-upper} no longer applies; this chapter does not
decide whether the strict-inequality encoding without that clause has an
upper limit of order \(\log^{1/3}n\).

\begin{hblemma}{thm:p03-w03-p03-071-log-mod}{proved=annex:p03-w03-p03-071-log-mod}{The residue of \(\lfloor\log u\rfloor\) is definable}
In the structure with universe \(S=\{1,\ldots,u\}\), its order and the
relations \(\mathrm{Add}_u\) and \(\mathrm{Mul}_u\), let
\begin{align*}
  \operatorname{One}(x)&:\ \neg\exists_y\,y<x,\\
  \operatorname{Ev}(a)&:\ \exists_b\,\mathrm{Add}(b,b,a),\\
  \operatorname{Pw}(w)&:\ \forall_a\forall_b\,(\mathrm{Mul}(a,b,w)\to
     \operatorname{One}(a)\vee\operatorname{Ev}(a)),\displaybreak[1]\\
  \operatorname{Top}(w)&:\ \operatorname{Pw}(w)\wedge\forall_{w'}
     (\operatorname{Pw}(w')\to w'\leq w),\\
  \delta_0(x)&:\ \operatorname{One}(x),\\
  \delta_{i+1}(x)&:\ \exists_y\,(\delta_i(y)\wedge\mathrm{Add}(y,y,x)),\displaybreak[1]\\
  \kappa(x)&:\ \exists_y\,(\delta_{40}(y)\wedge y=x+1),\\
  \operatorname{Cong}_r(w)&:\ \delta_r(w)\vee\exists_t\exists_v\exists_q
     \exists_h\,\bigl(\kappa(q)\wedge\delta_r(h)\\
  &\qquad\wedge\mathrm{Mul}(t,q,v)\wedge\mathrm{Add}(v,h,w)\bigr),\\
  \Phi&:\ \exists_w\,\Bigl(\operatorname{Top}(w)\wedge
     \bigvee_{r=1}^{20}\operatorname{Cong}_r(w)\Bigr).
\end{align*}
If \(u\geq2^{40}\), then \(\Phi\) holds in this structure if and only if
\(\lfloor\log u\rfloor\bmod 40\in\{1,\ldots,20\}\).
\end{hblemma}

Spencer says only that once \(\{1,\ldots,u\}\) is arithmetized ``we can talk
about \(\log u\)'' \cite[p.~23]{Spencer2007LogicRandomStructuresAuthor};
\cref{thm:p03-w03-p03-071-log-mod} supplies the formula.  The hypothesis
\(u\geq2^{40}\) is needed only so that the constant \(2^{40}-1\) lies in
\(S\).

Spencer also says that once \(\{1,\ldots,u\}\) is arithmetized one can say
that \(u\) is prime and that \(u\) is a Fermat prime
\cite[p.~23]{Spencer2007LogicRandomStructuresAuthor}.  The next lemma gives
the formulas, and the corollary after it transfers them to the random
ordered graph.

\begin{hblemma}{thm:p03-w03-p03-071-prime}{proved=annex:p03-w03-p03-071-prime}{Primes and Fermat primes are definable}
In the structure \((S,<,\mathrm{Add}_u,\mathrm{Mul}_u)\) with
\(S=\{1,\ldots,u\}\), with \(\operatorname{One}\) and \(\operatorname{Pw}\) as
in \cref{thm:p03-w03-p03-071-log-mod}, let
\begin{align*}
  \operatorname{Div}(a,b)&:\ \exists_c\,\mathrm{Mul}(a,c,b),\\
  \operatorname{Two}(w)&:\ \exists_y\,(\operatorname{One}(y)\wedge
     \mathrm{Add}(y,y,w)),\\
  \operatorname{Sq}(w)&:\ \exists_y\,\mathrm{Mul}(y,y,w),\displaybreak[1]\\
  \operatorname{Prime}(x)&:\ \neg\operatorname{One}(x)\wedge\forall_a\forall_b\,
     \bigl(\mathrm{Mul}(a,b,x)\\
  &\qquad\to\operatorname{One}(a)\vee\operatorname{One}(b)\bigr),\\
  \operatorname{Fer}(x)&:\ \operatorname{Prime}(x)\wedge\exists_v\exists_{v'}
     \bigl[x=v+1\wedge v=v'+1\wedge\operatorname{Pw}(v)\\
  &\qquad\wedge\forall_w\forall_{w'}\bigl((\operatorname{Pw}(w)\wedge w=w'+1
     \wedge\operatorname{Div}(w',v'))\\
  &\qquad\qquad\to\operatorname{Two}(w)\vee\operatorname{Sq}(w)\bigr)\bigr].
\end{align*}
For every \(x\in S\), \(\operatorname{Prime}(x)\) holds if and only if \(x\)
is prime, and \(\operatorname{Fer}(x)\) holds if and only if \(x\) is a
Fermat prime, that is, a prime of the form \(2^{2^m}+1\) with \(m\geq0\).
\end{hblemma}

The idea of \(\operatorname{Fer}\): with \(x=2^j+1\), the exponent \(j\) is a
power of \(2\) exactly when every divisor \(i\geq2\) of \(j\) is even, and
\(i\) divides \(j\) exactly when \(2^i-1\) divides \(2^j-1\), so divisibility
of exponents is read off from divisibility of numbers in \(S\) without
defining exponentiation.

\begin{hbcorollary}{thm:p03-w03-p03-071-prime-graph}{proved=annex:p03-w03-p03-071-prime-graph}{Primality of an encoded length is almost surely first-order}
Let \(\operatorname{Prime}^{\ast}(u)\) be the formula
\(\exists_c\exists_d\exists_{c'}\exists_{d'}\,[\operatorname{Good}(u,c,d,c',d')
\wedge\operatorname{Prime}^{P,T}_u(u)]\), where
\(\operatorname{Prime}^{P,T}_u\) is \(\operatorname{Prime}\) with quantifiers
relativized to \(w\leq u\), \(\mathrm{Add}\) replaced by \(P\) and
\(\mathrm{Mul}\) by \(T\), built from the same \(c,d,c',d'\) as in
\cref{def:p03-w03-p03-071-sentence}; define \(\operatorname{Fer}^{\ast}(u)\)
in the same way from \(\operatorname{Fer}\).  For every ordered graph and
every \(u\), \(\operatorname{Prime}^{\ast}(u)\) implies that \(u\) is prime
and \(\operatorname{Fer}^{\ast}(u)\) implies that \(u\) is a Fermat prime.
In \(G_<(n,\frac12)\), with probability tending to \(1\), for every \(u\)
with \(1\leq u\leq0.9L\), \(\operatorname{Prime}^{\ast}(u)\) holds if and only
if \(u\) is prime and \(\operatorname{Fer}^{\ast}(u)\) holds if and only if
\(u\) is a Fermat prime.
\end{hbcorollary}

\begin{hbdefinition}{def:p03-w03-p03-071-sentence}{The nonconvergent sentence \(A\)}
Spencer's sentence says: there exists \(u\) such that
\(\{1,\ldots,u\}\) is first order arithmetizable, \(\{1,\ldots,u+1\}\) is
not, and \(\log u\) modulo \(40\) is one of \(1,2,\ldots,20\)
\cite[p.~23]{Spencer2007LogicRandomStructuresAuthor}.  With the corrected
encoding it is the first-order sentence
\[
  A:\quad\exists_u\exists_{u'}\exists_c\exists_d\exists_{c'}\exists_{d'}\,
  \bigl[u'=u+1\wedge\operatorname{Good}(u,c,d,c',d')\wedge
  \neg\operatorname{Ar}(u')\wedge\Phi^{P,T}_u\bigr],
\]
where \(\Phi^{P,T}_u\) is \(\Phi\) of \cref{thm:p03-w03-p03-071-log-mod} with
quantifiers relativized to \(w\leq u\), \(\mathrm{Add}\) replaced by \(P\)
and \(\mathrm{Mul}\) by \(T\) (built from the same \(c,d,c',d'\)), and
\(\log u\) read as \(\lfloor\log u\rfloor\).
\end{hbdefinition}

\begin{hbtheorem}{thm:p03-w03-p03-071-chs}{proved=annex:p03-w03-p03-071-chs}{Nonconvergence (Compton, Henson and Shelah)}
For the sentence \(A\) of \cref{def:p03-w03-p03-071-sentence},
\[
  \liminf_{n\to\infty}\Pr[G_<(n,\tfrac12)\entails A]=0,\qquad
  \limsup_{n\to\infty}\Pr[G_<(n,\tfrac12)\entails A]=1 .
\]
In particular \(\lim_{n\to\infty}\Pr[G_<(n,\frac12)\entails A]\) does not
exist.
\end{hbtheorem}

Spencer attributes the theorem to Compton, Henson and Shelah; his text
spells the second name ``Hansen'' on pp.~22 and~25, while his reference list
has Henson \cite[pp.~22, 28]{Spencer2007LogicRandomStructuresAuthor}.  His
account of why the sentence works is that \(n\) fixes the largest
arithmetizable \(u\) up to a constant factor, so \(\log u\) up to an
additive constant, and the residue class then depends only on \(n\)
\cite[pp.~23--24]{Spencer2007LogicRandomStructuresAuthor}.  The proof in the
annex follows that account with the constants \(0.9\) and \(3\) of
\cref{thm:p03-w03-p03-071-representation,thm:p03-w03-p03-071-upper}.

\begin{hbexample}{def:p03-w03-p03-071-u-three}{Encoding \(\{1,2,3\}\)}
Take \(u=3\).  Then
\[
 \mathcal F_+(3)=\{\{1,2\},\{1,2,3\}\},
\]
from \(1+1=2\) and \(1+2=3\), and \(\mathcal F_\times(3)=\varnothing\), since
\(2\cdot2>3\).  Suppose \(c>3\), the trace of \(c\) on \(\{1,2,3\}\) is
\(\{1,2\}\) and that of \(c+1\) is \(\{1,2,3\}\); for fixed \(c\) this has
probability \(2^{-6}\) in \(G_<(n,\frac12)\).  The vertex \(c\) witnesses
exactly the triple \((1,1,2)\) of \(R^<_{c,c+2}\), because the strict
inequalities force \(z=2\) and \(x=y=1\); the vertex \(c+1\) witnesses
\((1,2,3)\) and \((2,1,3)\).  So \(R^<_{c,c+2}=\mathrm{Add}_3\).  With
\(c'=d'=4\), the relation \(T\) consists of the unit triples
\((1,1,1)\), \((2,1,2)\), \((3,1,3)\), \((1,2,2)\) and \((1,3,3)\),
which is \(\mathrm{Mul}_3\).
By \cref{thm:p03-w03-p03-071-witness} the formula
\(\operatorname{Good}(3,c,c+2,4,4)\) holds.  Under Spencer's literal
\(R_{c,c+2}\) the same vertex \(c\) would also witness \((1,2,2)\) and
\((2,1,2)\), which are not sums.
\end{hbexample}

\begin{hbwarning}{def:p03-w03-p03-071-speed}{How slowly the oscillation happens}
Along \(n\), the residue in \cref{def:p03-w03-p03-071-sentence} changes only
when \(\lfloor\log(0.8L)\rfloor\) changes, that is, when \((\log n)^{1/3}\)
doubles.  A run of values of \(n\) on which \(\Pr[A]\) is near \(1\) is
therefore astronomically long, and the theorem gives no information for
any particular \(n\): the proof needs \(0.9L-1\geq2^{40}\).  This is
consistent with \cref{thm:p03-w03-p03-071-shelah} below, which forbids rapid
oscillation.
\end{hbwarning}

\begin{hbdefinition}{def:p03-w03-p03-071-common-nbhd}{Common neighbourhoods in the sparse random graph}
In the random graph \(G(n,p)\) of \cref{def:p03-w03-p03-065-gnp}, on the
vertex set \(\{1,\ldots,n\}\) and with no built-in order, take
\(p=n^{-1/4}\).  For distinct vertices \(x_1,x_2,x_3,x_4\),
\(N(x_1,x_2,x_3,x_4)\) is the set of vertices other than \(x_1,\ldots,x_4\)
that are adjacent to all four \cite[p.~24]{Spencer2007LogicRandomStructuresAuthor}.
Since \(p^4=\frac1n\), a fixed vertex outside a fixed quadruple lies in its
common neighbourhood with probability \(\frac1n\).  Here \(\ln\) is the
natural logarithm, and \(\sigma(n)=\lfloor\ln n/(20\ln\ln n)\rfloor\) for
\(n\geq16\).
\end{hbdefinition}

\begin{hblemma}{thm:p03-w03-p03-071-all-sizes}{proved=annex:p03-w03-p03-071-all-sizes}{Common neighbourhoods of every small size}
In \(G(n,n^{-1/4})\), with probability tending to \(1\), for every integer
\(s\) with \(0\leq s\leq\sigma(n)\) there are distinct vertices
\(x_1,x_2,x_3,x_4\) with \(|N(x_1,x_2,x_3,x_4)|=s\).
\end{hblemma}

Spencer states that such sets exist of all sizes up to roughly
\(\ln n/\ln\ln n\) \cite[p.~24]{Spencer2007LogicRandomStructuresAuthor}.
\Cref{thm:p03-w03-p03-071-all-sizes} proves this with the factor
\(\frac1{20}\); the proof does not reach the constant \(1\), and no upper
limit on the sizes is proved here.  The factor is irrelevant to the use the
outline makes of the lemma, which needs sizes only up to a constant multiple
of \(\ln^{1/3}n\), and \(C\ln^{1/3}n\leq\sigma(n)\) for every constant \(C\)
once \(n\) is large.

\begin{proofboundary}{def:p03-w03-p03-071-sparse-outline}{Nonconvergence for \(G(n,n^{-1/4})\) without order: what the source outlines and what is not proved here}
Spencer also outlines the same method for the unordered random graph
\(G(n,n^{-1/4})\) \cite[p.~24]{Spencer2007LogicRandomStructuresAuthor}.
There a set \(S\) is arithmetized in second-order language by a binary
\(<\) and ternary \(+,\ast\).  For \(u\notin S\) the ternary relation
\(R_u(x,y,z)\) on \(S\) says that \(u,x,y,z\) have a common neighbour, and
for \(u,x\notin S\) the binary relation is \(R_{u,x}(y,z)=R_u(x,y,z)\); the
source notes that these relations are symmetric while \(<\) is not, and that
further work is needed for this.  \(S\) is first order arithmetizable if
there are \(u_1,u_2,u_3,u_4\) with \(R_{u_1,u_2},R_{u_3},R_{u_4}\) playing the
roles of \(<,+,\ast\).  Because one cannot quantify over all \(S\), the
encoding uses the common neighbourhoods \(S=N(x_1,x_2,x_3,x_4)\) of four
vertices (\cref{def:p03-w03-p03-071-common-nbhd}), compares sizes by \(\operatorname{BIGGER}(S,T)\), an injection
from \(T\) to \(S\) given by some \(R_{u_1,u_2}\), as in the almost sure
encoding of \cref{sec:p03-w03-p03-069-encoding}, and states that some
\(N(x_1,x_2,x_3,x_4)\) is arithmetizable while no bigger one is, and that
its size \(u\) has \(\log u\) between \(1\) and \(20\) modulo \(40\).  Two
steps are proved here: four given vertices have no common neighbour among
the other \(n-4\) with probability \((1-p^4)^{n-4}=(1-\frac1n)^{n-4}\), which
tends to \(e^{-1}\); and common neighbourhoods of every size up to
\(\sigma(n)\) occur almost surely (\cref{thm:p03-w03-p03-071-all-sizes}).
The source asserts without proof that each \(R_u\) behaves like an
independent random ternary relation of density \(1-e^{-1}\), that
\(+,\ast,<\) are almost surely represented for \(|S|\leq\ln^{1/3}n\), and
that a maximal arithmetizable \(N(x_1,x_2,x_3,x_4)\) has size
\(\Theta(\ln^{1/3}n)\); it also leaves open how the symmetric relations
\(R_{u,x}\) are to represent the asymmetric \(<\).  The same difficulty
affects \(+\) and \(\ast\): \(R_u\) is symmetric in all three arguments, so
taken literally it can never be addition on a set of at least two elements
(on \(\{1,\ldots,u\}\), \((1,1,2)\) is a sum and \((2,1,1)\) is not).  None of these claims is
proved in this chapter, and the resulting nonconvergence for
\(G(n,n^{-1/4})\) is not asserted here as a theorem.
\end{proofboundary}

\section{Slow oscillation and the evenness of \(n\)}
\label{sec:p03-w03-p03-071-evenness}

\begin{hbdefinition}{def:p03-w03-p03-071-coupling}{Restriction to a random subset, and almost sure representation}
Let \(N=2n+1\) and let \(G\) be an ordered graph on \(\{1,\ldots,N\}\).  For
\(S\subseteq\{1,\ldots,N\}\), \(G|S\) is the induced ordered graph on \(S\).
\(G^{(n)}\) is \(G|S\) for \(S\) a uniformly random subset of size exactly
\(n\), and \(G^{(n+1)}\) is \(G|S\) for \(S\) uniform of size \(n+1\); the
probabilities \(\Pr[G^{(k)}\entails A]\) are over \(S\) with \(G\) fixed.
\(\mu(G)\) is the probability that \(G_<(2n+1,p)\) is \(G\)
\cite[p.~25]{Spencer2007LogicRandomStructuresAuthor}.  A sentence \(A\)
\emph{almost surely represents evenness} if \(f_A(2n)\to1\) and
\(f_A(2n+1)\to0\) \cite[p.~25]{Spencer2007LogicRandomStructuresAuthor}.
\end{hbdefinition}

\begin{hblemma}{thm:p03-w03-p03-071-coupling}{proved=annex:p03-w03-p03-071-coupling}{The coupling identity}
For every \(k\leq N\) and every \(p\), if \(G\sim G_<(N,p)\) and \(S\) is an
independent uniform \(k\)-subset, then \(G|S\), transported to
\(\{1,\ldots,k\}\) by the order isomorphism, is distributed as
\(G_<(k,p)\).  Hence for every sentence \(A\)
\[
  f_A(n+1)-f_A(n)=\sum_G\mu(G)\Bigl[\Pr[G^{(n+1)}\entails A]
  -\Pr[G^{(n)}\entails A]\Bigr].
\]
\end{hblemma}

\begin{hbdefinition}{def:p03-w03-p03-071-circuit}{Circuits and total influence}
Following O'Donnell, for \(d\geq2\) a \emph{depth-\(d\) circuit} over
Boolean variables \(s_1,\ldots,s_N\) is a directed acyclic graph whose nodes
lie in \(d+1\) layers, with every arc from some layer \(j-1\) to layer
\(j\); layer \(0\) consists of exactly \(2N\) nodes labelled by the literals
\(s_b,\neg s_b\); layer \(d\) has exactly one node, the output; the nodes of
layers \(1,3,5,\ldots\) all carry one of the labels AND, OR and those of
layers \(2,4,\ldots\) the other; gates compute the AND or OR of their
in-neighbours \cite[def.~4.26, p.~103]{ODonnell2021AnalysisBooleanArxiv}.
Its \emph{size} is the number of nodes in layers \(1\) to \(d-1\), and no
layer-\(1\) node is joined to a variable or its negation more than once
\cite[def.~4.27, p.~104]{ODonnell2021AnalysisBooleanArxiv}.  For
\(f:\{0,1\}^N\to\{0,1\}\), coordinate \(b\) is pivotal at \(s\) if flipping
\(s_b\) changes \(f\); \(\operatorname{Inf}_b[f]\) is the probability that
\(b\) is pivotal at a uniformly random \(s\), and the total influence is
\(\mathbf I[f]=\sum_b\operatorname{Inf}_b[f]\)
\cite[defs.~2.12--2.13, p.~46; def.~2.27,
p.~49]{ODonnell2021AnalysisBooleanArxiv}.  O'Donnell writes truth values as
\(\pm1\); pivotality, and so \(\mathbf I\), does not depend on how the two
values are named.
\end{hbdefinition}

\begin{hbdefinition}{def:p03-w03-p03-071-restriction}{Restrictions, decision trees and normal forms}
Let \(V\) be a finite set of coordinates.  A \emph{restriction} of \(V\) is a
pair \((J\mid z)\) with \(J\subseteq V\) and \(z\in\{0,1\}^{V\setminus J}\);
the coordinates in \(J\) are \emph{free} and the others \emph{fixed}.  For
\(f:\{0,1\}^V\to\{0,1\}\), the restricted function
\(f_{J\mid z}:\{0,1\}^J\to\{0,1\}\) is \(f\) with the fixed coordinates set
to \(z\) \cite[def.~3.18, p.~75]{ODonnell2021AnalysisBooleanArxiv}.  For
\(0\leq\delta\leq1\), a \emph{\(\delta\)-random restriction} puts each
coordinate into \(J\) independently with probability \(\delta\) and then
chooses \(z\) uniformly; equivalently, each coordinate is independently free
with probability \(\delta\) and fixed to \(0\), or to \(1\), with probability
\(\frac{1-\delta}2\) each
\cite[def.~4.15, p.~98]{ODonnell2021AnalysisBooleanArxiv}.  If \((J'\mid
z')\) is a restriction of \(J\), the \emph{composed} restriction of \(V\) is
\((J'\mid(z,z'))\), and \((f_{J\mid z})_{J'\mid z'}=f_{J'\mid(z,z')}\).
Influences and total influence of a function on \(\{0,1\}^J\) are defined as
in \cref{def:p03-w03-p03-071-circuit} with \(J\) for the set of coordinates;
for \(J=\varnothing\) the total influence is \(0\).

A \emph{decision tree} over \(J\) is a rooted binary tree whose internal
nodes are labelled by coordinates in \(J\), whose two outgoing edges at each
internal node are labelled \(0\) and \(1\), whose leaves are labelled \(0\)
or \(1\), and in which no coordinate appears twice on a root-to-leaf path;
on input \(y\) it follows the edges labelled by the queried values of \(y\)
and outputs the label of the leaf reached.  Its depth is the largest number
of internal nodes on a root-to-leaf path, and \(\operatorname{DT}(g)\) is the
least depth of a decision tree computing \(g\)
\cite[defs.~3.13--3.14, pp.~73--74]{ODonnell2021AnalysisBooleanArxiv}; a
constant function has \(\operatorname{DT}=0\).  A \emph{term} is an AND, and
a \emph{clause} an OR, of literals on distinct coordinates; a \emph{DNF} is
an OR of terms, a \emph{CNF} an AND of clauses, and the \emph{width} is the
largest number of literals in a term or clause
\cite[def.~4.1, p.~93; def.~4.3, p.~94]{ODonnell2021AnalysisBooleanArxiv}.
An AND of no inputs is \(1\) and an OR of no inputs is \(0\), for terms and
clauses and for the gates of a circuit alike.
\end{hbdefinition}

\begin{hblemma}{thm:p03-w03-p03-071-restriction-facts}{proved=annex:p03-w03-p03-071-restriction-facts}{Three facts about restrictions}
Let \(V\) be a finite set, \(f:\{0,1\}^V\to\{0,1\}\) and
\(0\leq\delta,\delta_1,\delta_2\leq1\).
\begin{enumerate}
\item If \((J\mid z)\) is a \(\delta\)-random restriction of \(V\), then
  \(\mathbb E\bigl[\mathbf I[f_{J\mid z}]\bigr]=\delta\,\mathbf I[f]\).
\item If \((J\mid z)\) is a \(\delta_1\)-random restriction of \(V\) and,
  given it, \((J'\mid z')\) is a \(\delta_2\)-random restriction of \(J\),
  then the composed restriction \((J'\mid(z,z'))\) is a
  \(\delta_1\delta_2\)-random restriction of \(V\).
\item If \(g:\{0,1\}^J\to\{0,1\}\) has \(\operatorname{DT}(g)\leq k\), then
  \(\mathbf I[g]\leq k\), and \(g\) is computed by a DNF of width at most
  \(k\) and by a CNF of width at most \(k\).  Every \(g:\{0,1\}^J\to\{0,1\}\)
  has \(\mathbf I[g]\leq|J|\).
\end{enumerate}
\end{hblemma}

O'Donnell proves item~1 through Fourier coefficients
\cite[cor.~4.18, p.~99]{ODonnell2021AnalysisBooleanArxiv}, observes item~2
\cite[p.~104]{ODonnell2021AnalysisBooleanArxiv}, and proves the normal-form
part of item~3 \cite[prop.~4.5, p.~94]{ODonnell2021AnalysisBooleanArxiv};
the annex proves all three directly.

\begin{hblemma}{thm:p03-w03-p03-071-switching}{proved=annex:p03-w03-p03-071-switching}{Switching lemma (H\aa stad)}
Let \(F\) be a DNF or a CNF of width at most \(w\) over a finite set \(V\),
let \(0<\delta\leq\frac12\), and let \((J\mid z)\) be a \(\delta\)-random
restriction of \(V\).  Then for every integer \(k\geq1\),
\[
  \Pr\bigl[\operatorname{DT}(F_{J\mid z})\geq k\bigr]\leq(16\delta w)^k .
\]
\end{hblemma}

O'Donnell states the lemma with the constant \(5\) in place of \(16\) and
refers for its proof to H\aa stad and to Razborov
\cite[pp.~100--101]{ODonnell2021AnalysisBooleanArxiv}.  The annex gives a
complete proof by an encoding argument; the constant \(16\) is all that the
next theorem needs.

\begin{hbtheorem}{thm:p03-w03-p03-071-ac0-influence}{proved=annex:p03-w03-p03-071-ac0-influence}{Total influence of bounded-depth circuits}
Let \(N\geq1\), \(d\geq2\), and let \(f:\{0,1\}^N\to\{0,1\}\) be computed by
a depth-\(d\) circuit of size \(s\geq1\).  Put
\(t=\lceil3\log(2Ns)\rceil\).  Then
\[
  \mathbf I[f]\leq4\,(32t)^{d-1}\leq4\,\bigl(128\log(2Ns)\bigr)^{d-1}.
\]
\end{hbtheorem}

Boppana's theorem, as O'Donnell states it, is the sharper bound
\(\mathbf I[f]\leq O(\log s)^{d-1}\), with no dependence on \(N\)
\cite[thm.~4.30, p.~106]{ODonnell2021AnalysisBooleanArxiv}; O'Donnell
attributes it to Boppana \cite[p.~111]{ODonnell2021AnalysisBooleanArxiv} and
leaves its proof as Exercise~4.20
\cite[p.~106]{ODonnell2021AnalysisBooleanArxiv}.
\Cref{thm:p03-w03-p03-071-ac0-influence} has \(\log(2Ns)\) in place of
\(\log s\).  For circuits of size at most \(N^c\) the two agree up to a
factor depending on \(c\) and \(d\), because \(\log(2Ns)\leq(c+1)\log(2N)\);
that is the only case used below.  Whether the dependence on \(N\) can be
removed is not proved in this chapter.  The proof follows the plan of
O'Donnell's lemma~4.28 \cite[pp.~104--106]{ODonnell2021AnalysisBooleanArxiv}:
a \(\frac12\)-random restriction shortens the bottom gates, as in his
lemma~4.19 \cite[p.~99]{ODonnell2021AnalysisBooleanArxiv}; each further
restriction switches one layer of DNFs or CNFs into shallow decision trees,
so that two layers merge into one; and a final restriction turns the output
into a shallow decision tree, whose total influence is at most its depth.

The circuit fact Spencer names is different: a bounded-depth circuit of
polynomial size cannot compute majority, a theorem he credits originally to
Razborov, after which ``some further technical work'' shows that such a
circuit cannot tell a random set of \(n\) true inputs from one of \(n+1\)
\cite[p.~26]{Spencer2007LogicRandomStructuresAuthor}.  Failure to compute
majority exactly does not bound the distinguishing advantage, so the proof
below replaces that step by \cref{thm:p03-w03-p03-071-ac0-influence} and the
elementary \cref{thm:p03-w03-p03-071-slice}; their combination is
\cref{thm:p03-w03-p03-071-advantage}, the precise form of the
indistinguishability Spencer describes.  O'Donnell derives the majority
lower bound itself from Boppana's theorem
\cite[p.~107]{ODonnell2021AnalysisBooleanArxiv}.

\begin{hblemma}{thm:p03-w03-p03-071-translation}{proved=annex:p03-w03-p03-071-translation}{A first-order sentence becomes a bounded-depth circuit}
Let \(A\) be a sentence with \(\ell\) occurrences of connectives,
quantifiers and atomic formulas and quantifier rank \(q\), and let
\(d_A=\max(2,2\ell)\).  For every ordered graph \(G\) on
\(\{1,\ldots,N\}\), the function \(f_{G}:\{0,1\}^N\to\{0,1\}\) with
\(f_G(s)=1\) exactly when \(G|\{b:s_b=1\}\entails A\) is constant, or a
single literal, or computed by a depth-\(d\) circuit with \(2\leq d\leq d_A\)
and size at most \(s_N=(d_A+1)\,\ell\,(3N)^{q}+2\).
\end{hblemma}

\begin{hblemma}{thm:p03-w03-p03-071-slice}{proved=annex:p03-w03-p03-071-slice}{Adjacent slices}
For \(f:\{0,1\}^N\to\{0,1\}\) and \(0\leq k<N\), let \(\mu_k(f)\) be the
average of \(f\) over the inputs with exactly \(k\) ones.  Then
\[
  |\mu_{k+1}(f)-\mu_k(f)|\leq\frac{2^{N-1}\,\mathbf I[f]}{\binom Nk(N-k)} .
\]
For \(N=2n+1\) and \(k=n\) the right side is at most
\(\mathbf I[f]/\sqrt n\).
\end{hblemma}

\begin{hbcorollary}{thm:p03-w03-p03-071-advantage}{proved=annex:p03-w03-p03-071-advantage}{Bounded-depth circuits cannot tell \(n\) true inputs from \(n+1\)}
Let \(d\geq2\), \(c\geq1\) and \(n\geq1\).  If
\(f:\{0,1\}^{2n+1}\to\{0,1\}\) is computed by a depth-\(d\) circuit of size
\(s\) with \(1\leq s\leq(2n+1)^c\), then
\[
  |\mu_{n+1}(f)-\mu_n(f)|\leq
  \frac{4\,\bigl(128(c+1)\log(4n+2)\bigr)^{d-1}}{\sqrt n},
\]
which tends to \(0\) as \(n\to\infty\) for fixed \(d\) and \(c\).
\end{hbcorollary}

\begin{hbtheorem}{thm:p03-w03-p03-071-shelah}{proved=annex:p03-w03-p03-071-shelah}{Slow oscillation (Shelah)}
For every first-order sentence \(A\) there is a constant \(K\) such that, for
all \(n\geq2\) and every ordered graph \(G\) on \(2n+1\) vertices,
\[
  \bigl|\Pr[G^{(n+1)}\entails A]-\Pr[G^{(n)}\entails A]\bigr|
  \leq K\frac{(\log n)^{d_A-1}}{\sqrt n},
\]
with \(d_A\) as in \cref{thm:p03-w03-p03-071-translation}.  Consequently
\(\lim_{n\to\infty}\bigl(f_A(n+1)-f_A(n)\bigr)=0\) for \(G_<(n,\frac12)\).
\end{hbtheorem}

The theorem is Shelah's \cite[pp.~25, 28]{Spencer2007LogicRandomStructuresAuthor};
the uniform statement over all \(G\) is the form Spencer reports Shelah as
proving \cite[p.~25]{Spencer2007LogicRandomStructuresAuthor}.  The proof
uses the value of \(p\) only through \cref{thm:p03-w03-p03-071-coupling},
which holds for every \(p\), so the conclusion holds for every fixed \(p\).

\begin{hbcorollary}{thm:p03-w03-p03-071-evenness}{proved=annex:p03-w03-p03-071-evenness}{No almost sure representation of evenness}
No first-order sentence almost surely represents evenness in
\(G_<(n,\frac12)\): there is no first-order \(A\) with \(f_A(2n)\to1\) and
\(f_A(2n+1)\to0\).
\end{hbcorollary}

\begin{hbexample}{def:p03-w03-p03-071-translation-example}{Translating a sentence into a circuit}
Spencer illustrates the translation with a sentence and its relativization
to \(S\) \cite[p.~25]{Spencer2007LogicRandomStructuresAuthor}; the displayed
sentence and its displayed translation do not match, and the translation is
that of
\[
  B:\quad\exists_x\forall_y\exists_z\,(z\sim x\wedge z\sim y),
\]
``some vertex shares a neighbour with every vertex.''  (The displayed
sentence, with \(y\sim x\) inside \(\forall_y\), is false in every graph,
since \(y=x\) is allowed.)  Fix \(G\) on \(\{1,\ldots,N\}\) and write
\(N_G(b)\) for the neighbourhood of \(b\).  By the construction of
\cref{thm:p03-w03-p03-071-translation}, \(G|S\entails B\) exactly when
\[
  \bigvee_{b=1}^{N}\Bigl[s_b\wedge\bigwedge_{c=1}^{N}\Bigl(\neg s_c\vee
  \bigvee_{e\in N_G(b)\cap N_G(c)}s_e\Bigr)\Bigr]=1,
\]
the atomic facts \(e\sim b\), \(e\sim c\) having been evaluated in \(G\).
Layered, the inner disjunctions are \(N^2\) OR-gates at layer \(1\), the
literal \(s_b\) passes through one single-input OR-gate at layer \(1\), the
\(N\) conjunctions are AND-gates at layer \(2\), and the output OR is at
layer \(3\): a depth-\(3\) circuit of size \(N^2+2N\).  If some
\(N_G(b)\cap N_G(c)\) is empty the corresponding gate is
\(\neg s_c\) alone.  \Cref{thm:p03-w03-p03-071-ac0-influence,thm:p03-w03-p03-071-slice}
then give \(|\Pr[G^{(n+1)}\entails B]-\Pr[G^{(n)}\entails B]|\leq
4(128\log(2N^3+4N^2))^2/\sqrt n\) for every \(G\).
\end{hbexample}

\begin{hbtheorem}{thm:p03-w03-p03-071-mso}{proved=annex:p03-w03-p03-071-mso}{The first-order restriction is necessary}
Let \(\operatorname{EV}\) be the monadic second-order sentence
\[
  \exists X\,\bigl[\forall_x(\operatorname{One}(x)\to Xx)\wedge
  \forall_x\forall_y(y=x+1\to(Xx\leftrightarrow\neg Xy))\wedge
  \forall_x(\operatorname{Last}(x)\to\neg Xx)\bigr],
\]
where \(\operatorname{Last}(x)\) is \(\neg\exists_y\,x<y\).  An ordered graph
on \(n\geq1\) vertices satisfies \(\operatorname{EV}\) if and only if \(n\)
is even.  Hence \(f_{\operatorname{EV}}(2n)=1\) and
\(f_{\operatorname{EV}}(2n+1)=0\) for every \(p\), and
\cref{thm:p03-w03-p03-071-shelah,thm:p03-w03-p03-071-evenness} fail for
monadic second-order sentences.
\end{hbtheorem}

\begin{hbwarning}{def:p03-w03-p03-071-not-convergence}{Slow oscillation is not convergence}
\Cref{thm:p03-w03-p03-071-shelah} does not say that \(f_A(n)\) converges:
the sentence of \cref{thm:p03-w03-p03-071-chs} has \(f_A(n+1)-f_A(n)\to0\)
and still no limit.  Nor does it bound how far \(f_A(n)\) and \(f_A(n+t)\)
are apart when \(t\) grows with \(n\).  The order is essential to the
nonconvergence: without it, \(G(n,\frac12)\) satisfies the zero--one law
(\cref{thm:p03-w03-p03-067-zero-one}).
\end{hbwarning}

\section{Relationships}
\label{sec:p03-w03-p03-071-relationships}

\Cref{thm:p03-w03-p03-071-evenness} implies that parity of a finite linear
order is not first-order definable in the order alone.  If a sentence
\(A\) in \(<\) alone held in exactly the orders of even size, then read in
an ordered graph, where it ignores \(\sim\), it would have
\(f_A(2n)=1\) and \(f_A(2n+1)=0\), contrary to the corollary.  The direct
proof by Ehrenfeucht games is \cref{def:p03-ef-parity}; the probabilistic
statement is stronger, since \(A\) may use \(\sim\) and need only be right
with probability tending to \(1\).  The unordered constant-probability
theory is in \cref{sec:p03-w03-p03-067-constant-p}, and the almost sure
encoding of a second-order property by common neighbourhoods, which the
outline of \cref{def:p03-w03-p03-071-sparse-outline} reuses, is in
\cref{sec:p03-w03-p03-069-encoding}.  Libkin's remark that even cardinality
has no zero--one truth limit is cited in \cref{sec:p09-presentation-n-c-002}.
Spencer's reference list sends the reader to Compton, Henson and Shelah for
results on the random ordered graph and to Shelah for the evenness theorem
\cite[p.~28]{Spencer2007LogicRandomStructuresAuthor}; those papers were not
consulted for this chapter.

\section{Sources}
\label{sec:p03-w03-p03-071-sources}

Spencer, ``Logic and Random Structures'', author's version, sections 12 and
13, pp.~22--26, with ``About the references'', pp.~27--28
\cite{Spencer2007LogicRandomStructuresAuthor}.  O'Donnell, \emph{Analysis
of Boolean Functions}, arXiv edition of May 2021, definitions 2.12--2.13,
2.27, 3.13--3.14, 3.18, 4.1, 4.3, 4.15, 4.26--4.27, proposition 4.5,
corollary 4.18, lemmas 4.19 and 4.28, the switching lemma (pp.~100--101) and
theorem 4.30 with exercise 4.20 (p.~106) \cite{ODonnell2021AnalysisBooleanArxiv}.
